\documentclass[EJP,preprint]{ejpecp}
% Supplement to the paper problem131_switching.tex (Paper C of Problem 131).
% It builds on its own. The statement numbers of the paper are read from
% problem131_switching.aux, so build the paper first for them to resolve.
\usepackage{booktabs}
\usepackage[shortlabels]{enumitem}
\usepackage{xr}
\externaldocument[M-]{problem131_switching}
\newcommand{\mref}[1]{\ref*{M-#1}}
\newcommand{\meqref}[1]{(\mref{#1})}
\renewcommand{\thesection}{S\arabic{section}}

\newcommand{\Z}{\mathbb{Z}}
\newcommand{\R}{\mathbb{R}}
\newcommand{\E}{\mathbb{E}}
\newcommand{\Prob}{\mathbb{P}}
\newcommand{\supp}{\operatorname{supp}}
\newcommand{\ones}{\mathbf{1}}
\newcommand{\tree}{\operatorname{tree}}

\TITLE{Supplement to ``Face selection for slowly switching Markov chains''}
\SHORTTITLE{Supplement to ``Face selection for slowly switching Markov chains''}
\AUTHORS{Arjun Pemmasani\footnote{Department of Mathematics, Harvey Mudd College,
Claremont, CA 91711, USA.\\ \EMAIL{apemmasani@hmc.edu}}}
\KEYWORDS{occupation times of Markov chains; mode; principal Dirichlet eigenvalue}
\AMSSUBJ{60J10; 60J27; 05C50}
\ABSTRACT{This supplement holds proofs for the paper ``Face selection for slowly
switching Markov chains''. Section~\ref{supp:examples} proves the rows of the table of
outcomes (Table~\mref{tab:outcomes} of the paper) that are proved by hand and not in the text,
namely the path $P_d$ and two triangles joined by a link. Section~\ref{supp:threecycle} gives
the full proof of Proposition~\mref{prop:threecycle} of the paper, the directed 3-cycle, whose
proof the paper only sketches. Sections~\ref{supp:laplace-next} and~\ref{supp:nextterm}
give the higher-order Laplace proof and the detailed next-term expansion, including
the complete-graph coefficient calculation. Numbers with the prefix S refer to this supplement, and all other
numbers refer to the paper.}

\begin{document}

\section{Proofs for the table of outcomes}\label{supp:examples}\label{app:examples}

This section proves the rows $P_d$ and $\Theta_\varepsilon$ of
Table~\mref{tab:outcomes}. In both rows $Q$ is symmetric and the start has full
support. So the admissible faces are the connected faces
(Lemma~\mref{lem:admissible}), and the winners at $\tau$ are the connected faces
that minimize $\Phi_\tau(F)=\tau\lambda_F+\lvert F\rvert-1$
(Theorem~\mref{thm:hull}). We use three facts about a connected face $F$. First,
$-Q_F$ is symmetric, $\lambda_F$ is its least eigenvalue, and an eigenvector of
$-Q_F$ with positive entries belongs to $\lambda_F$ (the second fact at the
start of Appendix~\mref{app:torus} of the paper). Second, $-Q_F$ is the Laplacian of the
graph on $F$, which is positive semidefinite, plus the diagonal matrix of the
rates out of $F$. So $\lambda_F$ is the least value of
\begin{equation}\label{eq:rayleigh}
 \sum_{\{i,j\}\subseteq F}q_{ij}(f_i-f_j)^2+\sum_{i\in F}q_i^{\rm out}(F)f_i^2
\end{equation}
over unit vectors $f\in\R^F$. Third, if $F'\subsetneq F$ is nonempty, then
$\lambda_{F'}>\lambda_F$ (Proposition~\mref{prop:rate}(d)).

Each proof has two steps, namely the exit rates and then the lower-left boundary
of the hull. For the second step we use the following lemma. It says that a
chain of points with increasing switch values is its own lower-left boundary.

\begin{lemma}[a chain of lines]\label{lem:chain}
Let $x_0>x_1>\dots>x_m\ge0$ and $y_0<y_1<\dots<y_m$, and put
$\ell_i(\tau)=\tau x_i+y_i$ and $t_i=(y_i-y_{i-1})/(x_{i-1}-x_i)$ for
$1\le i\le m$. Put $t_0=0$ and $t_{m+1}=\infty$, and suppose that
$t_1<t_2<\dots<t_m$. Then for $1\le j\le m$ and every $\tau$
\begin{equation}\label{eq:chain}
 \ell_j(\tau)-\ell_{j-1}(\tau)=(x_{j-1}-x_j)(t_j-\tau).
\end{equation}
Let $0\le i\le m$ and let $\tau$ be finite with $t_i\le\tau\le t_{i+1}$. Then
$\ell_0(\tau)\ge\dots\ge\ell_i(\tau)\le\dots\le\ell_m(\tau)$. If also
$t_i<\tau<t_{i+1}$, then all these inequalities are strict, so
$\ell_i(\tau)<\ell_j(\tau)$ for every $j\ne i$.
\end{lemma}

\begin{proof}
The identity~\eqref{eq:chain} is the definition of $t_j$. If $j\le i$, then
$t_j\le t_i\le\tau$, so $\ell_j(\tau)\le\ell_{j-1}(\tau)$. If $j>i$, then
$\tau\le t_{i+1}\le t_j$, so $\ell_j(\tau)\ge\ell_{j-1}(\tau)$. If
$t_i<\tau<t_{i+1}$, these inequalities are strict.
\end{proof}

\subsection{The path}\label{app:paths}

Let $P_d$ be the path with states $1,\dots,d$ and $q_{i,i+1}=q_{i+1,i}=1$ for
$1\le i<d$. An \emph{end-arc} of $k<d$ states is one of the two faces
$\{1,\dots,k\}$ and $\{d-k+1,\dots,d\}$. For $k\ge1$ put
\[
 \lambda'_k=2-2\cos\frac{\pi}{2k+1},\qquad
 \tau'_k=\frac1{\lambda'_{k-1}-\lambda'_k}\ (k\ge2),\qquad
 D'_k=k+\lambda'_k\tau'_k\ (k\ge2).
\]

\begin{lemma}[the numbers $D'_k$]\label{lem:Dprime}
\begin{enumerate}[(a)]
\item The sequence $\lambda'_k$ is positive, strictly decreasing and strictly
convex in $k\ge1$. Hence $0<\tau'_2<\tau'_3<\cdots$. Also
$\tau'_2=(1+\sqrt5)/2$.
\item For every $k\ge2$, $\frac32k-\frac12<D'_k<\frac32k$.
\end{enumerate}
\end{lemma}

\begin{proof}
(a) Let $g(x)=2-2\cos(\pi/x)$, so that $\lambda'_k=g(2k+1)$. For $x\ge2$,
\[
 g'(x)=-\frac{2\pi\sin(\pi/x)}{x^2}<0,\qquad
 g''(x)=2\pi\Bigl(\frac{\pi\cos(\pi/x)}{x^4}+\frac{2\sin(\pi/x)}{x^3}\Bigr)>0 .
\]
So $\lambda'_{k-1}-\lambda'_k>\lambda'_k-\lambda'_{k+1}>0$ for $k\ge2$, and
the $\tau'_k$ increase. Since $\lambda'_1=1$ and
$\lambda'_2=2-2\cos\frac\pi5=(3-\sqrt5)/2$, we get
$\tau'_2=2/(\sqrt5-1)=(1+\sqrt5)/2$.

(b) We turn the claim into two bounds on the ratio of two consecutive exit
rates, and we prove each bound by comparing $\sin x$ with $x$. Put
$a=\pi/(2(2k-1))$ and $b=\pi/(2(2k+1))$, so that $0<b<a\le\pi/6$. Since
$2-2\cos2x=4\sin^2x$, we have $\lambda'_{k-1}=4\sin^2a$ and
$\lambda'_k=4\sin^2b$. So the ratio $\rho=\lambda'_k/\lambda'_{k-1}$ equals
$(\sin b/\sin a)^2$. By definition
$D'_k-k=\lambda'_k/(\lambda'_{k-1}-\lambda'_k)=\rho/(1-\rho)$, which increases
with $\rho$ and takes the values $\frac{k-1}2$ and $\frac k2$ at
$\rho=\frac{k-1}{k+1}$ and $\rho=\frac k{k+2}$. So the claim is
\[
 \frac{k-1}{k+1}<\rho<\frac k{k+2}.
\]
If we replace each sine by its argument, then $\rho$ becomes
$(b/a)^2=(2k-1)^2/(2k+1)^2$. This number lies between the two bounds and very
close to the lower one. So the lower bound needs only $\sin b/\sin a>b/a$, and
the upper bound needs the next term of the sine.

\emph{The lower bound.} The function $\sin x/x$ decreases on $(0,\pi]$, so
$\sin b/\sin a>b/a=(2k-1)/(2k+1)$. Then $\rho>(2k-1)^2/(2k+1)^2>(k-1)/(k+1)$,
where the last step holds because $(2k-1)^2(k+1)-(k-1)(2k+1)^2=2>0$.

\emph{The upper bound.} Write $s(x)=\sin x/x$ and
$y=(8k-2)/((k+2)(2k-1)^2)$. Then
\[
 \rho=\Bigl(\frac{2k-1}{2k+1}\Bigr)^2\Bigl(\frac{s(b)}{s(a)}\Bigr)^2
 \qquad\text{and}\qquad
 \frac k{k+2}=\Bigl(\frac{2k-1}{2k+1}\Bigr)^2(1+y),
\]
the second because $k(2k+1)^2-(k+2)(2k-1)^2=8k-2$. So we must show that
$s(b)/s(a)<\sqrt{1+y}$. We reduce this in three steps to an inequality between
two explicit functions of $k$.

Step 1. We remove the square root. We have $0<y<1$, because
$(k+2)(2k-1)^2\ge12(2k-1)>8k-2$, and $\sqrt{1+y}\ge1+\frac y3$ for
$0\le y\le3$. So it is enough to show that $s(b)/s(a)<1+\frac y3$, that is
\[
 s(b)-s(a)<\tfrac y3\,s(a).
\]

Step 2. We bound both sides with
$1-\frac{x^2}6\le s(x)\le1-\frac{x^2}6+\frac{x^4}{120}$, which holds for
$x>0$. On the right, $a\le\pi/6$ gives $s(a)\ge1-\frac{\pi^2}{216}>0.954$, so
$\frac y3s(a)>0.318\,y$. On the left,
\[
 s(b)-s(a)\le\frac{a^2-b^2}6+\frac{b^4}{120}
 =\frac{\pi^2k}{3(4k^2-1)^2}+\frac{\pi^4}{1920(2k+1)^4}.
\]
So it is enough to show that this is less than $0.318\,y$.

Step 3. We multiply both sides by $(2k-1)^2$, and on the left we use
$k/(2k+1)^2<1/(4k)$ and $(2k-1)^2/(2k+1)^4<1/(4k^2)$. It is then enough to show
that
\[
 \frac{\pi^2}{12k}+\frac{\pi^4}{7680k^2}<0.318\,\frac{8k-2}{k+2}.
\]
The left side decreases in $k$, and the right side, which is
$0.318(8-\frac{18}{k+2})$, increases. At $k=2$ the left side is less than
$0.415$ and the right side is more than $1.11$. So the inequality holds for
every $k\ge2$.
\end{proof}

\begin{proposition}[the path]\label{prop:paths}
Let $d\ge3$, let $Q$ be the path $P_d$ with unit rates, and let the start have
full support. Put $K=\lfloor2d/3\rfloor$ and
$\tau_{\rm full}=\lceil d/3\rceil/\lambda'_K$.
\begin{enumerate}[(a)]
\item The connected faces are the sets of consecutive states. An end-arc of
$k<d$ states has exit rate $\lambda'_k$. A set of $k$ consecutive states that
contains neither $1$ nor $d$ has exit rate $2-2\cos\frac\pi{k+1}>\lambda'_k$.
\item We have $2\le K\le d-1$ and
$\tau'_2<\dots<\tau'_K<\tau_{\rm full}$.
\item The winners are the states $1$ and $d$ for $0<\tau<\tau'_2$, the two
end-arcs of $k$ states for $\tau'_k<\tau<\tau'_{k+1}$ and $2\le k\le K-1$, the
two end-arcs of $K$ states for $\tau'_K<\tau<\tau_{\rm full}$, and $P_d$ alone
for $\tau>\tau_{\rm full}$. No other face wins for any $\tau>0$.
\end{enumerate}
\end{proposition}

In words, the mode grows from an end of the path, one state at a time, until
it covers $K=\lfloor2d/3\rfloor$ states (about two thirds of the path), and then it jumps to the whole path. The
first step is at the golden ratio.

\emph{Idea of the proof.} Growing an end-arc by one state lowers its exit rate
by less and less, since $\lambda'_k$ is convex. So the end-arcs alone win in
the order of their size, and the arc of $k$ states starts to win at
$\tau'_k$. The whole path enters when its exponent $d-1$ falls below the least
exponent of the arcs. At $\tau'_k$ that least exponent is $D'_k-1$, so the
arcs of $k$ states win before the whole path exactly when $D'_k<d$.
Lemma~\ref{lem:Dprime}(b) puts $D'_k$ just below $\frac32k$, and this gives
$k\le2d/3$.

\begin{proof}
(a) A connected face of a path is a set of consecutive states. Let $F$ have
$k<d$ states. If $F=\{1,\dots,k\}$, then $-Q_F$ is tridiagonal with
off-diagonal entries $-1$ and diagonal $(1,2,\dots,2)$. Put
$\vartheta=\pi/(2k+1)$ and $r_j=\cos((2j-1)\vartheta/2)$ for $1\le j\le k$.
Then $r_j>0$, since $(2j-1)\vartheta/2<\pi/2$. With $r_0=r_1$ and
$r_{k+1}=\cos(\pi/2)=0$ we have $r_{j-1}+r_{j+1}=2\cos\vartheta\,r_j$ for
$1\le j\le k$, and this says that $-Q_Fr=(2-2\cos\vartheta)r$. So
$\lambda_F=\lambda'_k$, and the same holds for $\{d-k+1,\dots,d\}$ by the
reflection $i\mapsto d+1-i$. If $F$ contains neither $1$ nor $d$, then
$-Q_F=2I-A$ with $A$ the adjacency matrix of a path on $k$ states. Now
$r_j=\sin(j\pi/(k+1))$ is a positive eigenvector with eigenvalue
$2-2\cos\frac\pi{k+1}$, and this is larger than $\lambda'_k$ because
$\frac\pi{k+1}>\frac\pi{2k+1}$.

(b) Since $d\ge3$ we have $2\le K\le d-1$. The $\tau'_k$ increase by
Lemma~\ref{lem:Dprime}(a). Note that $d-K=\lceil d/3\rceil$, so
$\tau_{\rm full}=(d-K)/\lambda'_K$. By Lemma~\ref{lem:Dprime}(b),
$D'_K<\frac32K\le d$, that is $\lambda'_K\tau'_K<d-K$. So
$\tau'_K<\tau_{\rm full}$.

(c) Write $\phi_k(\tau)=\tau\lambda'_k+k-1$ for the exponent of an end-arc of
$k$ states, $1\le k\le d-1$. First we apply Lemma~\ref{lem:chain} to the
points $(\lambda'_k,k-1)$ for $1\le k\le K$ followed by $(0,d-1)$, the point
of $P_d$. Its switch values are $\tau'_2,\dots,\tau'_K$ and
$(d-K)/\lambda'_K=\tau_{\rm full}$, which increase by (b). So on each of the
intervals in (c) the face named there has a smaller exponent than the other
end-arcs of at most $K$ states and than $P_d$, and by (a) also than every set
of at most $K$ consecutive states that is not an end-arc.

It remains to show that a face $F$ of $k$ consecutive states with
$K<k<d$ never wins. By (a) its exponent is at least $\phi_k(\tau)$. We show that
$\phi_k(\tau)>\min(\phi_K(\tau),d-1)$ for every $\tau>0$. Now we apply
Lemma~\ref{lem:chain} to the points $(\lambda'_j,j-1)$ for $1\le j\le d-1$,
with switch values $\tau'_2<\dots<\tau'_{d-1}$. If $\tau<\tau'_{K+1}$, then
$\tau<\tau'_j$ for $K<j\le k$, and~\eqref{eq:chain} gives
$\phi_k(\tau)>\phi_{k-1}(\tau)>\dots>\phi_K(\tau)$. If $\tau\ge\tau'_{K+1}$, take $j$
with $K+1\le j\le d-1$ and $\tau'_j\le\tau$, and with $\tau\le\tau'_{j+1}$ if
$j<d-1$. Then
\[
 \phi_k(\tau)\ge \phi_j(\tau)\ge \phi_j(\tau'_j)=D'_j-1>\tfrac32j-\tfrac32
 \ge\tfrac32K\ge d-1 .
\]
The first step is Lemma~\ref{lem:chain}, the second holds since $\phi_j$
increases, the third is the definition of $D'_j$, the fourth is
Lemma~\ref{lem:Dprime}(b), the fifth is $j\ge K+1$, and the last holds because
$3K\ge2d-2$. So $F$ never wins. The winners at $\tau$ are the admissible faces with the
least exponent, and this gives (c).
\end{proof}

At a switch value the faces on both sides tie. On $P_4$ we have $K=2$ and
$D'_3=4.077>4$, so the end-arc of three states never wins. Its point lies
only $\lambda'_3-\frac12\lambda'_2=0.0071$ to the right of the edge from the
end-arcs of two states to $P_4$. The distance is small because the lower bound
in Lemma~\ref{lem:Dprime}(b) is nearly attained.

\subsection{Two triangles}\label{app:triangles}

Let $\varepsilon>0$. The graph $\Theta_\varepsilon$ has states $1,\dots,6$,
unit rates on the triangles $A=\{1,2,3\}$ and $B=\{4,5,6\}$, and one more link
between $3$ and $4$ of rate $\varepsilon$. We call $1,2,5,6$ the inner states
and $\{1,2\}$, $\{5,6\}$ the inner edges. The graph has the symmetries
$1\leftrightarrow2$, $5\leftrightarrow6$ and the mirror map $i\mapsto7-i$. Put
\begin{align*}
 p_A(x)&=x^2-(3+\varepsilon)x+\varepsilon,\\
 p_4(x)&=x^3-(5+2\varepsilon)x^2+(6+6\varepsilon)x-2\varepsilon,\\
 p_5(x)&=x^3-(4+2\varepsilon)x^2+(3+4\varepsilon)x-\varepsilon,
\end{align*}
and $\varepsilon_0=(64-\sqrt{46})/60=0.95363$,
$\varepsilon_1=(64+\sqrt{46})/60=1.17971$, the roots of
$120\varepsilon^2-256\varepsilon+135$.

\begin{proposition}[two triangles joined by a link]\label{prop:triangles}
Let the start have full support.
\begin{enumerate}[(a)]
\item The two triangles have the same exit rate
$\lambda_A=\lambda_B=\frac12\bigl(3+\varepsilon-\sqrt{\varepsilon^2+2\varepsilon+9}\bigr)$,
which is the smaller root of $p_A$, and $0<\lambda_A<1$.
\item Among the connected faces of one, two and three states the least exit
rates are $2$, $1$ and $\lambda_A$. They are attained exactly by the inner
states, by the inner edges and by $A$ and $B$. Put
$\lambda_4=\lambda_{\{1,2,3,4\}}$ and $\lambda_5=\lambda_{\{2,\dots,6\}}$. Every
connected face of four states has exit rate at least $\lambda_4$, and every
connected face of five states has exit rate $\lambda_5$. Here $\lambda_4$ and
$\lambda_5$ are the least roots of $p_4$ and $p_5$, and
$\lambda_5<\lambda_4<\lambda_A$.
\item For every $\varepsilon>0$, $\lambda_5>\frac12\lambda_4$, and
$\lambda_4>\frac23\lambda_A$ if and only if $\varepsilon<\varepsilon_1$.
\item Let $\varepsilon<\varepsilon_1$, and put
\[
 \tau_A=\frac1{1-\lambda_A}=\frac{1+\varepsilon+\sqrt{\varepsilon^2+2\varepsilon+9}}4,
 \qquad
 \tau^*=\frac3{\lambda_A}=\frac{3\bigl(3+\varepsilon+\sqrt{\varepsilon^2+2\varepsilon+9}\bigr)}{2\varepsilon}.
\]
Then $1<\tau_A<\tau^*$. The winners are the inner states for $0<\tau<1$, the
inner edges for $1<\tau<\tau_A$, the triangles $A$ and $B$ for
$\tau_A<\tau<\tau^*$, and $[6]$ alone for $\tau>\tau^*$. No other face wins
for any $\tau>0$. As $\varepsilon\to0$,
$\tau_A=1+\frac\varepsilon3+O(\varepsilon^2)$ and
$\tau^*=\frac9\varepsilon+2+O(\varepsilon)$.
\item Let $\varepsilon>\varepsilon_1$. Then at $\tau=3/\lambda_A$, where the
triangles and $[6]$ have the same exponent, the face $\{1,2,3,4\}$ has a
smaller exponent. So the list in (d) fails.
\end{enumerate}
\end{proposition}

In words, for a weak link each triangle is a community that the chain rarely
leaves. The mode sits on a triangle from $\tau\approx1$ until
$\tau\approx9/\varepsilon$, and then it spreads over all six states at once.
Once the link is stronger than $\varepsilon_1$ this list fails, because a
triangle together with the far end of the link beats both the triangles and
$[6]$ at the value of $\tau$ where they tie.

\emph{Idea of the proof.} The exit rates of small faces are explicit, and
those of the faces with four and five states are roots of cubics. The points
of the inner states, the inner edges, the triangles and $[6]$ form a chain as
in Lemma~\ref{lem:chain}. The points of the faces with four and five states
must lie strictly to the right of the edge from a triangle to $[6]$, and this
is what (c) says. We prove (c) by continuity in
$\varepsilon$. A change of sign needs a common root of two polynomials, hence a
zero of a resultant, and the resultants are explicit polynomials in
$\varepsilon$.

\begin{proof}
(a) On the vectors $(a,a,b)$ the matrix $-Q_A$ acts as
$(a,b)\mapsto(a-b,\,-2a+(2+\varepsilon)b)$, with characteristic polynomial
$p_A$, and $(1,-1,0)$ is an eigenvector with eigenvalue $3$. As
$p_A(0)=\varepsilon>0$ and $p_A(1)=-2<0$, the smaller root of $p_A$ lies in
$(0,1)$ and the other root is larger than $1$. So the smaller root is
$\lambda_A$. The mirror map gives $\lambda_B=\lambda_A$.

(b) By~\eqref{eq:rayleigh}, $\lambda_F\ge1$ whenever every state of $F$ has
rate at least $1$ out of $F$. We go through the connected faces by size. The
states $3$ and $4$ have exit rate $2+\varepsilon$, and the inner states have
$2$. The connected pairs are the inner edges, with $\lambda=1$, the pair
$\{3,4\}$, with $\lambda=2$, and $\{1,3\}$, $\{2,3\}$ and their mirror images,
with $\lambda=\frac12(4+\varepsilon-\sqrt{\varepsilon^2+4})>1$. The connected
faces of three states are $A$, $B$ and the sets $\{a,3,4\}$ and $\{3,4,b\}$
with $a\in\{1,2\}$ and $b\in\{5,6\}$, because a connected face that meets both
triangles contains $3$ and $4$. In each of these four sets every state has
rate at least $1$ out of the set, so $\lambda\ge1>\lambda_A$. A connected face
of four states meets both triangles, so it contains $3$ and $4$. So it is
$\{1,2,3,4\}$, its mirror image $\{3,4,5,6\}$, or a set $\{a,3,4,b\}$, in which every state has rate $1$ out
of the set, so that $\lambda\ge1$. A connected face of five states is the
complement of an inner state (the complements of $3$ and of $4$ are
disconnected), and these four faces are images of $\{2,\dots,6\}$
under the symmetries. By the third fact above,
$\lambda_5<\lambda_{\{3,4,5,6\}}=\lambda_4<\lambda_A<1$, since
$\{2,\dots,6\}\supsetneq\{3,4,5,6\}$ and $\{1,2,3,4\}\supsetneq A$.

For the polynomials, note that $(1,-1,0,0)$ is an eigenvector of $-Q_F$ for
$F=\{1,2,3,4\}$ with eigenvalue $3$. On the vectors $(a,a,b,c)$ the matrix
acts as
$(a,b,c)\mapsto(a-b,\,-2a+(2+\varepsilon)b-\varepsilon c,\,-\varepsilon b+(2+\varepsilon)c)$,
with characteristic polynomial
$(x-1)(x-2)(x-2-2\varepsilon)-2(x-2-\varepsilon)=p_4(x)$. So the eigenvalues
of $-Q_F$ are $3$ and the roots of $p_4$. As $\lambda_4<1$, it is the least
root of $p_4$. Now let $F=\{2,\dots,6\}$. The vector that is $1$ at $5$, $-1$
at $6$ and $0$ elsewhere is an eigenvector of $-Q_F$ with eigenvalue $3$. On
the vectors with the value $a$ at $2$, $b$ at $3$, $c$ at $4$ and $e$ at both
$5$ and $6$, the matrix acts as
\[
 (a,b,c,e)\mapsto\bigl(2a-b,\ -a+(2+\varepsilon)b-\varepsilon c,\
 -\varepsilon b+(2+\varepsilon)c-2e,\ -c+e\bigr).
\]
Expanding the determinant along the first row
gives the characteristic polynomial
\[
 (x-2)\bigl[(x-1)\bigl((x-2-\varepsilon)^2-\varepsilon^2\bigr)-2(x-2-\varepsilon)\bigr]
 -\bigl[(x-2-\varepsilon)(x-1)-2\bigr]=(x-3)\,p_5(x).
\]
So the eigenvalues of $-Q_F$ are $3$ (twice) and the roots of $p_5$. As
$\lambda_5<1$, it is the least root of $p_5$. (Both characteristic polynomials are
also checked by computer algebra, see Section~\ref{app:examples-checks}.)

(c) Put $D_1=\lambda_5-\frac12\lambda_4$ and
$D_2=\lambda_4-\frac23\lambda_A$. Both are continuous in $\varepsilon>0$,
because the least eigenvalue of a symmetric matrix is continuous in its
entries. If $D_1=0$, then $x=\lambda_5$ is a common root of $p_5(x)$ and
$p_4(2x)$. If $D_2=0$, then $x=\lambda_4$ is a common root of $p_4(x)$ and
$p_A(\frac32x)$. A common root makes the resultant vanish, and an exact
computation gives
\begin{align*}
 \operatorname{Res}_x\bigl(p_5(x),p_4(2x)\bigr)
 &=-8\varepsilon^3(16\varepsilon^2+12\varepsilon+9),\\
 \operatorname{Res}_x\bigl(p_A(\tfrac32x),p_4(x)\bigr)
 &=-\tfrac1{16}\varepsilon^2(120\varepsilon^2-256\varepsilon+135).
\end{align*}
So $D_1$ has no zero in $\varepsilon>0$, and $D_2$ can vanish only at
$\varepsilon_0$ and at $\varepsilon_1$. Hence $D_1$ has one sign on
$(0,\infty)$, and $D_2$ has one sign on each of the intervals
$(0,\varepsilon_0)$, $(\varepsilon_0,\varepsilon_1)$ and
$(\varepsilon_1,\infty)$. A zero of the resultant does not force $D_2$ to
vanish, because the common root need not be the least root of $p_4$. So we find
the sign of $D_2$ for small $\varepsilon$, for large $\varepsilon$ and at
$\varepsilon_0$.

First let $\varepsilon$ be small. By (a) and (b),
$0<\lambda_5<\lambda_4<\lambda_A<\varepsilon$. Here $\lambda_5>0$ because
$\lambda_5\ge0$ by~\eqref{eq:rayleigh} and $p_5(0)=-\varepsilon\ne0$, and
$\lambda_A<\varepsilon$ because $p_A(\varepsilon)=-2\varepsilon<0$. At $\varepsilon=0$ the cubics are
$x(x-2)(x-3)$ and $x(x-1)(x-3)$, with the simple root $0$. So the implicit
function theorem gives one root $\frac\varepsilon3+O(\varepsilon^2)$ of each
cubic near $0$, and no other root near $0$. So
$\lambda_4=\frac\varepsilon3+O(\varepsilon^2)$ and
$\lambda_5=\frac\varepsilon3+O(\varepsilon^2)$, and (a) gives
$\lambda_A=\frac\varepsilon3+O(\varepsilon^2)$. Hence
$D_1=\frac\varepsilon6+O(\varepsilon^2)$ and
$D_2=\frac\varepsilon9+O(\varepsilon^2)$ are positive for small
$\varepsilon$. So $D_1>0$ for every $\varepsilon>0$, and $D_2>0$ on
$(0,\varepsilon_0)$.

Next, $D_2<0$ for large $\varepsilon$. Indeed $\lambda_A\to1$ as
$\varepsilon\to\infty$ by (a). On the other hand take $f=(a,a,c,c)$
in~\eqref{eq:rayleigh} for $\{1,2,3,4\}$. The value is $2(a-c)^2+2c^2$ and
$\lvert f\rvert^2=2a^2+2c^2$, so
$\lambda_4\le\min_{a,c}((a-c)^2+c^2)/(a^2+c^2)=(3-\sqrt5)/2<\frac12$ for every
$\varepsilon$.

Now we show that $D_2>0$ at $\varepsilon_0$. Put $x_0=\frac7{40}$. Then
\[
 p_4(x_0)=\frac{57743}{64000}-\frac{809}{800}\varepsilon,\qquad
 p_A(\tfrac32x_0)=\frac{59}{80}\varepsilon-\frac{4599}{6400}
\]
are both negative for $0.8923<\varepsilon<0.9743$, and this interval contains
$\varepsilon_0$. For $0\le x\le\frac12$ we have
$p_4'(x)=3x^2-2(5+2\varepsilon)x+6+6\varepsilon\ge1+4\varepsilon>0$. So $p_4$
increases on $[0,x_0]$ and is negative there, and its least root
$\lambda_4$ is larger than $x_0$. Also $p_A(\frac32x_0)<0$ puts $\frac32x_0$
between the roots of $p_A$, so $\lambda_A<\frac32x_0$. Hence
$\lambda_4>x_0>\frac23\lambda_A$ at $\varepsilon_0$.

So $D_2>0$ on $(0,\varepsilon_0]$. It has no zero in
$(\varepsilon_0,\varepsilon_1)$, so $D_2>0$ on $(0,\varepsilon_1)$. It has no
zero in $(\varepsilon_1,\infty)$ and is negative for large $\varepsilon$, so
$D_2<0$ on $(\varepsilon_1,\infty)$. By continuity $D_2=0$ at
$\varepsilon_1$. This proves (c).

(d) As $p_A(\frac34)=\frac\varepsilon4-\frac{27}{16}<0$ for
$\varepsilon<\frac{27}4$, we have $\lambda_A<\frac34$. So $1-\lambda_A>0$ and
$\lambda_A<3(1-\lambda_A)$, that is $1<\tau_A<\tau^*$. Apply
Lemma~\ref{lem:chain} to the points $(2,0)$, $(1,1)$, $(\lambda_A,2)$ and
$(0,5)$ of the inner states, the inner edges, the triangles and $[6]$. The
switch values are $1$, $\tau_A$ and $\tau^*$. So on each interval in (d) the
faces named there have a smaller exponent than the other three kinds, and by (b)
than every other connected face of at most three states. Let $F$ be a
connected face of four or five states. By (b) and (c), for every $\tau>0$
\begin{align*}
 \Phi_\tau(F)&\ge\tau\lambda_4+3>\tfrac23(\tau\lambda_A+2)+\tfrac13\cdot5
 &&\text{if }\lvert F\rvert=4,\\
 \Phi_\tau(F)&=\tau\lambda_5+4>\tfrac13(\tau\lambda_A+2)+\tfrac23\cdot5
 &&\text{if }\lvert F\rvert=5,
\end{align*}
where the second line uses $\lambda_5>\frac12\lambda_4>\frac13\lambda_A$. Each
right side is at least the smaller of the exponents $\tau\lambda_A+2$ of a
triangle and $5$ of $[6]$. So $F$ never wins, and this gives the winners. The closed forms follow from (a) by rationalizing, since
$(\varepsilon^2+2\varepsilon+9)-(1+\varepsilon)^2=8$ and
$(3+\varepsilon)^2-(\varepsilon^2+2\varepsilon+9)=4\varepsilon$. The
expansions follow from
$\sqrt{\varepsilon^2+2\varepsilon+9}=3+\frac\varepsilon3+O(\varepsilon^2)$.

(e) At $\tau=3/\lambda_A$ the triangles and $[6]$ have the exponent $5$, and
$\{1,2,3,4\}$ has $3\lambda_4/\lambda_A+3<5$ by (c).
\end{proof}

At $\varepsilon=1$, for example, $\tau_A=(1+\sqrt3)/2$ and
$\tau^*=6+3\sqrt3$.

\subsection{What was checked by computation}\label{app:examples-checks}

The proofs above use computer algebra in one place only, namely the two
resultants in the proof of Proposition~\ref{prop:triangles}(c). The script
\texttt{verify\_examples.py} computes them in sympy. It also expands the
characteristic polynomials of $-Q_F$ for $F=A$, $\{1,2,3,4\}$ and
$\{2,\dots,6\}$ and finds $(x-3)p_A$, $(x-3)p_4$ and $(x-3)^2p_5$, and it
recomputes $p_4(\frac7{40})$ and $p_A(\frac{21}{80})$ in rational arithmetic.
The rest of its checks for these two rows are in floating point, and no proof
rests on them. For $3\le d\le9$ it computes the exit rates of all connected
faces of $P_d$, compares those of the end-arcs with $\lambda'_k$, and finds
the sizes of the winners as in Proposition~\ref{prop:paths}(c). With 50
digits it confirms $\frac32k-\frac12<D'_k<\frac32k$ for $2\le k\le2000$. For
$\varepsilon=0.3$, $1$ and $1.1$ it computes the exit rates of all 63 faces of
$\Theta_\varepsilon$ and finds the least exit rates, the winners and the
switch values of Proposition~\ref{prop:triangles}. For $\varepsilon=1.5$, $3$,
$4.4$, $4.6$ and $8$ it finds that the faces $\{1,2,3,4\}$ and $\{3,4,5,6\}$
win on an interval of $\tau$, which agrees with
Proposition~\ref{prop:triangles}(e).

\section{The directed 3-cycle}\label{supp:threecycle}

This section proves Proposition~\mref{prop:threecycle} of the paper, whose statement we
recall. Let $d=3$, $q_{12}=b>2$, $q_{23}=q_{31}=1$, all other rates $0$, and $\mu_1=1$.
\begin{enumerate}[(i)]
\item The admissible faces are $\{1\}$, $\{1,2\}$ (reducible, $\lambda=1$) and
$\{1,2,3\}$ (irreducible, $\lambda=0$). Exactly
$M_{\{1\}}=(1-bh)^{N-1}$ and $M_{\{1,2\}}=bh(1-h)^{N-2}$, and the latter is
attained only at $n=(1,N-1,0)$. For the full face,
$\alpha^*=(1,b,b)/(1+2b)$, $\det\nabla^2I=(1+2b)^4/(3b^2)$ and
$K_{\{1,2,3\}}=(1+2b)^2/(2\sqrt3\,\pi b)$.
\item The equation $M_{\{1\}}=M_{\{1,2\}}$ has exactly one root $T^{(1)}_N$ with $bh<1$. It
satisfies $(b-1)T+\log T=L-\log b+O(T^2/N)$, so
$T^{(1)}_N=\frac1{b-1}[L-\log L+\log\frac{b-1}b]+O(\log L/L)$.
\item Every root $T\in[L/2,2L]$ of $M_{\{1,2\}}=M_{\{1,2,3\}}$ satisfies
$T=L+\log\frac{2\sqrt3\,\pi b^2}{(1+2b)^2}+O(1/L)$, and these roots lie in an interval of
length $O(1/L)$.
\item For large $N$ and $T\le N^{1/3}$,
the mode is $(N,0,0)$ for $T<T^{(1)}_N$, it is $(1,N-1,0)$ between
$T^{(1)}_N$ and the smallest root in (iii), and it lies on the full face above
the largest root in (iii).
\end{enumerate}

\begin{proof}
(i) \emph{The faces.} From state $1$ the chain can only move along $1\to2\to3\to1$. So a path
that starts at $1$ has support $\{1\}$, $\{1,2\}$ or $\{1,2,3\}$, and these are the admissible
faces. Inside $\{1,2\}$ the chain cannot return from $2$ to $1$, so this face is reducible with
components $\{1\}$ and $\{2\}$, of exit rates $b$ and $1$, and $\lambda_{\{1,2\}}=1$. The only
path with support $\{1\}$ stays at $1$ for $N-1$ steps, which gives $M_{\{1\}}$. A path with
support $\{1,2\}$ is $1^{n_1}2^{n_2}$, of probability $(1-bh)^{n_1-1}\,bh\,(1-h)^{n_2-1}$.
Raising $n_1$ by one multiplies this by $(1-bh)/(1-h)<1$, so the maximum is at $n_1=1$ only,
and it is $bh(1-h)^{N-2}$.

\emph{The full face.} Here $Q$ is a generator on $\{1,2,3\}$, so $\lambda=0$ and $r=\ones$. The
left vector $l$ is the stationary law. In it the three flows $b\,l_1$, $l_2$ and $l_3$ around the
cycle are equal, so $l=\alpha^*=(1,b,b)/(1+2b)$ and $(\mu\cdot r)(l\cdot\ones)=1$. By
Proposition~\mref{prop:constant}(a) it remains to find $\det\nabla^2I$ at $\alpha^*$, and then
$K_{\{1,2,3\}}=\sqrt{\det\nabla^2I}/(2\pi)$.

We first find the rate in closed form. By definition
\[
I(\alpha)=\sum_iq_i\alpha_i-\inf_{u>0}\Bigl[b\alpha_1\frac{u_2}{u_1}+\alpha_2\frac{u_3}{u_2}
+\alpha_3\frac{u_1}{u_3}\Bigr].
\]
The three ratios are positive numbers with product $1$, and every such triple occurs. So by the
AM--GM inequality the infimum is $3(b\alpha_1\alpha_2\alpha_3)^{1/3}$, and
\[
I(\alpha)=\sum_iq_i\alpha_i-3b^{1/3}g(\alpha),\qquad g(\alpha)=(\alpha_1\alpha_2\alpha_3)^{1/3}.
\]
The first sum is linear, so the Hessian of $I$ is $-3b^{1/3}$ times the Hessian of $g$. In all
three coordinates, $\partial_i\partial_jg=-g\,\Omega_{ij}$ with
$\Omega=\frac13\operatorname{diag}(w)^2-\frac19ww^\top$ and $w=1/\alpha$ (entrywise). In the
coordinates $(\alpha_2,\alpha_3)$ we have $\alpha_1=1-\alpha_2-\alpha_3$, so raising $\alpha_2$ or
$\alpha_3$ moves $\alpha$ in the direction $v_2=e_2-e_1$ or $v_3=e_3-e_1$. Hence $\nabla^2I$ is
the $2\times2$ matrix with entries $3b^{1/3}g\,v_i^\top\Omega v_j$.

Now we evaluate at $\alpha^*$. Put $s=1+2b$ and $y=1/b$. Then $3b^{1/3}g=3b/s$ and
$w=s(1,y,y)$, and a short computation gives
\[
v_2^\top\Omega v_2=v_3^\top\Omega v_3=\tfrac29s^2(1+y+y^2),\qquad
v_2^\top\Omega v_3=\tfrac19s^2(2+2y-y^2).
\]
The determinant of this $2\times2$ matrix is a difference of two squares, namely
\[
\frac{s^4}{81}\bigl[(2+2y+2y^2)^2-(2+2y-y^2)^2\bigr]=\frac{s^4}{81}\cdot3y^2(2+y)^2
=\frac{s^6}{27b^4},
\]
since $2+y=s/b$. Hence $\det\nabla^2I=(3b/s)^2s^6/(27b^4)=s^4/(3b^2)$ and
$K_{\{1,2,3\}}=s^2/(2\sqrt3\,\pi b)$. The exact parts of (i) are checked in Lean.

(ii) We show that $D_1=\log M_{\{1,2\}}-\log M_{\{1\}}$ increases in $h$. By (i),
$D_1=\log(bh)+(N-2)\log(1-h)-(N-1)\log(1-bh)$, so
\[
D_1'(h)=\frac1h-\frac{N-2}{1-h}+\frac{b(N-1)}{1-bh}>0\qquad(bh<1),
\]
because $b>1$ gives $b/(1-bh)>1/(1-h)$. Also $D_1\to-\infty$ as $h\to0$ and $D_1\to\infty$ as
$bh\to1$. So there is exactly one root (this is checked in Lean). To locate it we expand in
$T=Nh$, which gives
\[
D_1=\log b+\log T-L+(b-1)T+O(T^2/N).
\]
This is negative at $T=L/(2(b-1))$ and positive at $T=L/(b-1)$ for large $N$. So the root lies
between them, and it satisfies the equation in (ii). Since the root is of order $L$, that
equation gives $\log T=\log L-\log(b-1)+O(\log L/L)$, and putting this back gives the expansion.

(iii) By (i), $\log M_{\{1,2\}}=\log b+\log T-L-T+O(T^2/N)$. The full face satisfies (H), and
Theorem~\mref{thm:facemax} with $k=3$ and $\lambda=0$ gives
$\log M_{\{1,2,3\}}=-2L+\log T+\log K_{\{1,2,3\}}+O(1/T+T^2/N)$. So on $[L/2,2L]$
\[
\log M_{\{1,2\}}-\log M_{\{1,2,3\}}=L-T+\log b-\log K_{\{1,2,3\}}+O(1/L).
\]
At a root the left side is $0$, which gives the formula in (iii) with the value of
$K_{\{1,2,3\}}$ from (i). Two roots then differ by $O(1/L)$.

(iv) We compare the three face maxima on three ranges of $T$. Note first that
Lemma~\mref{lem:upper} with $\theta=q$ gives $M_{\{1,2,3\}}\le C(T/N)^2$ for $T\le N^{1/3}$.

Let $T<T^{(1)}_N$. Then $M_{\{1\}}>M_{\{1,2\}}$ by (ii). Also $T<L/(b-1)$ by (ii), so
$M_{\{1\}}=(1-bh)^{N-1}\ge N^{-b/(b-1)-o(1)}$. This is larger than $C(T/N)^2$ because
$b>2$ gives $b/(b-1)<2$. So the mode is $(N,0,0)$.

Let $T>T^{(1)}_N$. Then $M_{\{1,2\}}>M_{\{1\}}$ by (ii), and we compare $\{1,2\}$ with the full
face. For $T\le L/2$ the bound on $M_{\{1,2,3\}}$ gives
$\log M_{\{1,2\}}-\log M_{\{1,2,3\}}\ge L-T-\log T-C>0$. On $[L/2,2L]$ the
difference in (iii) is continuous in $T$, positive at $L/2$ and negative at
$2L$. So it is positive below its smallest root and negative above its largest root. For
$2L\le T\le N^{1/3}$ the same expansions give a difference of at most $-L+C<0$. Since the
maximum of $M_{\{1,2\}}$ is attained only at $(1,N-1,0)$, this proves (iv).
\end{proof}




\section{The higher-order Laplace estimate}\label{supp:laplace-next}

This section proves Lemma~\mref{lem:laplace2} of the paper. Its hypotheses,
coefficient formula and notation are stated in Appendix~\mref{app:laplace-next} there.

In words, $c_1$ is the mean of an explicit polynomial of degree $6$ under the
Gaussian law with covariance $\mathcal J^{-1}$. As a formal check of the
coefficients in~\meqref{eq:laplace-c1}, take the integrand
$e^{T(e^{\mathrm i\varphi}-1-\mathrm i\varphi)}$ with $m=1$, $a=1$ and integer
$T$. Its integral over $[-\pi,\pi]$ is $2\pi e^{-T}T^T/T!$, and Stirling's
formula gives $(2\pi/T)^{1/2}(1-\frac1{12T}+O(T^{-2}))$. Here $\mathcal J=1$,
$\Psi_{111}=-\mathrm i$ and $\Psi_{1111}=1$, so \meqref{eq:laplace-c1} gives
$c_1=\frac18-\frac1{12}-\frac18=-\frac1{12}$ as well. The check is only
formal, because $\Psi=e^{\mathrm i\varphi}-1-\mathrm i\varphi$ is not
$2\pi$-periodic. So it is not a function on the torus, and the lemma does not
cover it as stated. Only $e^{T\Psi}$ is periodic, for integer $T$.

\emph{Idea of the proof.} After the substitution $\varphi=x/\sqrt T$ the
integrand is a Gaussian weight times a power series in $T^{-1/2}$ whose
coefficients are polynomials in $x$. The coefficients of the odd powers of
$T^{-1/2}$ are odd polynomials and integrate to $0$. The coefficient of
$T^{-1}$ has the Gaussian mean $c_1$, by Wick's formula.

\begin{proof}
We follow the proof of Lemma~\mref{lem:laplace} with its notation, and expand
further. Now $B_\Psi$ and $B_a$ are the maxima over $|\beta|\le6$ and
$|\beta|\le4$. The localization and the scaling are as before. The integral
over $\lVert\varphi\rVert\ge\varepsilon$ is at most
$(2\pi)^mB_ae^{T(\Psi(\alpha,0)-\eta)}\le CT^{-2-m/2}e^{T\Psi(\alpha,0)}$.

\emph{Taylor expansion.} Put $s=T^{-1/2}$, so $0<s\le1$. For $j=3,4,5$ let
$\mathcal P_j(x)$ be the part of degree $j$ of the Taylor polynomial of
$\Psi(\alpha,\cdot)$ at $0$. So $\mathcal P_3=P$ and
$\mathcal P_4(x)=\frac1{24}\Psi_{ijkl}x^ix^jx^kx^l$. For $j=1,2,3$ let $A_j(x)$ be the
part of degree $j$ of the Taylor polynomial of $a(\alpha,\cdot)$, so
$A_1(x)=a_ix^i$ and $A_2(x)=\frac12a_{ij}x^ix^j$. Taylor's theorem to order
$6$ for $\Psi$ and to order $4$ for $a$ gives, on the cube
$\lVert x\rVert\le\varepsilon\sqrt T$,
\[
w=w_0+\omega_1,\qquad w_0=s\mathcal P_3+s^2\mathcal P_4+s^3\mathcal P_5,\qquad
a(\alpha,sx)=a_0+sA_1+s^2A_2+s^3A_3+\omega_2,
\]
with $|\omega_1|\le Cs^4|x|^6$ and $|\omega_2|\le Cs^4|x|^4$. Let
$\mathsf p(v)=1+v+\frac12v^2+\frac16v^3$. We compare $e^w$ with
$\mathsf p(w)$, then with $\mathsf p(w_0)$, and then with
\[
Y=1+s\mathcal P_3+s^2\bigl(\mathcal P_4+\tfrac12\mathcal P_3^2\bigr)
+s^3\bigl(\mathcal P_5+\mathcal P_3\mathcal P_4+\tfrac16\mathcal P_3^3\bigr),
\]
which collects the terms of $\mathsf p(w_0)$ of degree at most $3$ in $s$.
First, $|e^v-\mathsf p(v)|\le\frac1{24}|v|^4e^{|v|}$, and we know that
$|w|\le Cs|x|^3$ and $|w|\le\frac14\sigma_{\min}|x|^2$. So
$|e^w-\mathsf p(w)|\le Cs^4|x|^{12}e^{\sigma_{\min}|x|^2/4}$. Second,
$\mathsf p(w)-\mathsf p(w_0)
=\omega_1\bigl[1+\frac12(w+w_0)+\frac16(w^2+ww_0+w_0^2)\bigr]$, and $|w|$ and
$|w_0|$ are at most $C(1+|x|)^5$, by $|w|\le Cs|x|^3$ and $s\le1$. So
$|\mathsf p(w)-\mathsf p(w_0)|\le Cs^4(1+|x|)^{16}$. Third,
$\mathsf p(w_0)-Y$ is a sum of terms $s^j$ with $4\le j\le9$ times
polynomials in $x$ of degree at most $15$ with bounded coefficients, so
$|\mathsf p(w_0)-Y|\le Cs^4(1+|x|)^{15}$. Together,
\[
|e^w-Y|\le Cs^4(1+|x|)^{16}e^{\sigma_{\min}|x|^2/4}.
\]
Now write $e^wa=Y(a_0+sA_1+s^2A_2+s^3A_3)+Y\omega_2+(e^w-Y)a$. Here
$|Y|\le C(1+|x|)^9$ and $|a|\le B_a$. In the first product the terms of
degree at least $4$ in $s$ are at most $Cs^4(1+|x|)^{12}$. Hence
\[
e^{w}a(\alpha,sx)=a_0+s\,\mathcal Q_1(x)+s^2\mathcal Q_2(x)+s^3\mathcal Q_3(x)
+\omega(x),\qquad
|\omega(x)|\le CT^{-2}(1+|x|)^{16}e^{\sigma_{\min}|x|^2/4},
\]
where
\begin{align*}
\mathcal Q_1&=a_0\mathcal P_3+A_1,\qquad
\mathcal Q_2=a_0\bigl(\mathcal P_4+\tfrac12\mathcal P_3^2\bigr)+A_1\mathcal P_3+A_2,\\
\mathcal Q_3&=a_0\bigl(\mathcal P_5+\mathcal P_3\mathcal P_4+\tfrac16\mathcal P_3^3\bigr)
+A_1\bigl(\mathcal P_4+\tfrac12\mathcal P_3^2\bigr)+A_2\mathcal P_3+A_3 .
\end{align*}
Every term of $\mathcal Q_1$ and of $\mathcal Q_3$ has odd degree, so
$\mathcal Q_1$ and $\mathcal Q_3$ are odd functions.

\emph{Integration.} As before, $\omega$ contributes at most $CT^{-2}$ to the
integral over the cube, since $(1+|x|)^{16}e^{-\sigma_{\min}|x|^2/4}$ is
integrable. Extending the integrals of the four polynomial terms to $\R^m$
changes them by at most $Ce^{-\sigma_{\min}\varepsilon^2T/4}\le CT^{-2}$. On $\R^m$
the odd terms $\mathcal Q_1$ and $\mathcal Q_3$ integrate to $0$ against
$e^{-\frac12x^\top\mathcal Jx}$. This is why the error is $O(T^{-2})$ and not
$O(T^{-3/2})$. So the integral over the cube is
\[
e^{T\Psi(\alpha,0)}T^{-m/2}\Bigl[(2\pi)^{m/2}\det\mathcal J^{-1/2}
\Bigl(a_0+\frac{\E\mathcal Q_2(x)}T\Bigr)+O(T^{-2})\Bigr],
\]
where $x$ is a centered Gaussian vector with $\E[x^ix^j]=\mathcal J^{ij}$. As
in the proof of Lemma~\mref{lem:laplace}, the errors are $O(T^{-2})$ relative
to the main factor, because $\det\mathcal J\le(mB_\Psi)^m$.

It remains to show that $\E\mathcal Q_2=c_1$. By Wick's formula a Gaussian
moment of order $4$ or $6$ is the sum, over the three or fifteen ways to pair the
indices, of the products of the covariances of the pairs. The derivatives are
symmetric in their indices. So $\E[A_2]=\frac12a_{ij}\mathcal J^{ij}$,
$\E[A_1\mathcal P_3]=\frac16\cdot3\,a_i\Psi_{jkl}\mathcal J^{ij}\mathcal J^{kl}$ and
$\E[\mathcal P_4]=\frac1{24}\cdot3\,\Psi_{ijkl}\mathcal J^{ij}\mathcal J^{kl}$. In
$\E[\mathcal P_3^2]=\frac1{36}\Psi_{ijk}\Psi_{lrs}\E[x^ix^jx^kx^lx^rx^s]$, six pairings
join each index of the first factor to an index of the second, and each gives
$\Psi_{ijk}\Psi_{lrs}\mathcal J^{il}\mathcal J^{jr}\mathcal J^{ks}$. The other
nine pair two indices inside each factor and join the remaining two, and each
gives $\Psi_{ijk}\Psi_{lrs}\mathcal J^{ij}\mathcal J^{kl}\mathcal J^{rs}$. So
$\frac12\E[\mathcal P_3^2]$ gives the last two terms of~\meqref{eq:laplace-c1}. Finally
$c_1$ is a polynomial in $a_0$, in the derivatives at $\varphi=0$ and in the
entries of $\mathcal J^{-1}$, which are continuous on $\mathcal K$.
\end{proof}

\section{Proof of the next-term expansion}\label{supp:nextterm}

This section proves Proposition~\mref{prop:nextterm} of the paper. The coefficient
$c(\alpha)$ is defined in \meqref{eq:c-alpha}, using \meqref{eq:laplace-c1}.
All hypotheses and notation are those of the proposition and Appendix~\mref{app:nextterm}.

Throughout, $F$ satisfies (H). If $F=\{i\}$, then
\meqref{eq:vertexmax} gives $\log M_F=\log\mu_i-q_iT+O(T^2/N)$, which
is~\meqref{eq:logmaxnext} with $c_F=0$. So let $k\ge2$ until the proof of
(iii).

\emph{Part \textup{(i)}.} We use the proof of Theorem~\mref{thm:local} in
Appendix~\mref{app:torus}. It shows that $\Psi$ and $\chi$ in~\meqref{eq:A-fT}
satisfy the hypotheses of Lemma~\mref{lem:laplace} with $m=k-1$, and that their
$\varphi$-derivatives of every order are jointly continuous on
$\mathcal K\times\mathbb T^{k-1}$. So Lemma~\mref{lem:laplace2} applies with
$a=\chi$. Let $c_1(\alpha)$ be the function~\meqref{eq:laplace-c1} for this
$\Psi$ and $a=\chi$, with
$\mathcal J(\alpha)=-\nabla^2_\varphi\Psi(\alpha,0)$, and use the definition \meqref{eq:c-alpha} of the paper.
Then the bracket $\chi(\alpha,0)+O(1/T)$ in the last step of that proof
becomes $\chi(\alpha,0)+c_1(\alpha)/T+O(T^{-2})$. Since $\chi(\alpha,0)$ is
bounded below on $\mathcal K$, this is
$\chi(\alpha,0)(1+c(\alpha)/T+O(T^{-2}))$, and the rest of that proof gives
the expansion in (i). By the regularity step of that proof, $\Psi$ and $\chi$
are real-analytic in $(\alpha,\varphi)$ near
$\{\alpha_0\}\times\mathbb T^{k-1}$ for every $\alpha_0\in\Delta_F^\circ$. So
their $\varphi$-derivatives at $\varphi=0$ are real-analytic in $\alpha$, and
so are the entries of $\mathcal J(\alpha)^{-1}$ and $c$. Finally (i) gives
$c(\alpha)=\lim_{T\to\infty}T\bigl[f_T(\alpha)T^{-(k-1)/2}C_F(\alpha)^{-1}
e^{TI_F(\alpha)}-1\bigr]$. So $c$ is real, and it does not depend on which
state of $F$ we number $1$.

\emph{Part \textup{(ii)}, the plan.} We write $\log P_N(n)$ as a fixed term plus a function
$G_T$ of $\alpha=n/N$, and we maximize $G_T$. Far from $\alpha^*_F$, at distance more than a
large multiple of $1/T$, the quadratic growth of $I_F$ makes $G_T$ smaller than $(c_F-1)/T$.
At distance at most that multiple of $1/T$, the function $G_T$ is a constant plus a linear term
minus a positive definite quadratic form, up to an error $O(T^{-2})$. Completing the square shows
that this quadratic model has the maximum $c_F/T$, and that it takes it exactly at the point
$\alpha'=\alpha^{*\prime}_F+\Sigma_F\gamma_F/T$. So $G_T\le c_F/T+O(T^{-2})$, which gives the upper
bound. A lattice point next to that point gives the lower bound.

\emph{Part \textup{(ii)}, the constants.} Write $\Sigma=\Sigma_F$,
$\gamma=\gamma_F$ and $c^*=c(\alpha^*_F)$, and put
\[
X=N^{-(k-1)}T^{(k-1)/2}K_Fe^{-T\lambda_F},\qquad x=\alpha'-\alpha^{*\prime}_F .
\] By Propositions~\mref{prop:rate}(a)
and~\mref{prop:constant}(a) and by (i), the functions $I_F$, $\log C_F$ and
$c$ are real-analytic near $\alpha^*_F$, and $I_F$ has its minimum
$\lambda_F$ at $\alpha^*_F$ with Hessian $\Sigma^{-1}$. So there are a closed
ball $U=\{|x|\le\rho_0\}\subset\Delta_F^\circ$ and constants $c_U>0$ and
$L_1$ such that on $U$
\begin{gather*}
\Bigl|\log\frac{C_F(\alpha)}{K_F}-\gamma\cdot x\Bigr|\le L_1|x|^2,\qquad
\Bigl|I_F(\alpha)-\lambda_F-\tfrac12x^\top\Sigma^{-1}x\Bigr|\le L_1|x|^3,\\
|c(\alpha)-c^*|\le L_1|x|,\qquad|c(\alpha)|\le L_1,\qquad
I_F(\alpha)-\lambda_F\ge c_U|x|^2 .
\end{gather*}
Put $L_2=|\gamma|+L_1\rho_0$, so that $|\log(C_F(\alpha)/K_F)|\le L_2|x|$ on
$U$. Fix $\rho_1$ so large that
\[
\rho_1\ge\frac{L_2}{c_U},\qquad\rho_1\ge|\Sigma\gamma|+k,\qquad
c_U\rho_1^2-L_2\rho_1-L_1\ge|c_F|+1 .
\]
For $\alpha\in U$ let
\[
G_T(\alpha)=\log\frac{C_F(\alpha)}{K_F}-T\bigl(I_F(\alpha)-\lambda_F\bigr)
+\frac{c(\alpha)}T .
\]
By (i) with $\mathcal K=U$ and Lemma~\mref{lem:discrete}, there are $T_0$ and
$N_0$ such that for $N\ge N_0$, $T_0\le T\le N^{1/3}$ and $\supp n=F$ with
$\alpha=n/N\in U$,
\begin{equation}\label{eq:GT}
\log\frac{P_N(n)}X=G_T(\alpha)+O\Bigl(\frac1{T^2}+\frac{T^2}N\Bigr).
\end{equation}
Here we used $\log(1+c/T+O(T^{-2}))=c/T+O(T^{-2})$, which holds because
$|c|\le L_1$ on $U$. We take $T_0\ge\rho_1/\rho_0$, and we increase $T_0$ and
$N_0$ finitely many times below.

\emph{Part \textup{(ii)}, two bounds for $G_T$.} First let $\alpha\in U$ with
$|x|\ge\rho_1/T$. Then
$G_T(\alpha)\le L_2|x|-c_UT|x|^2+L_1/T$. The function
$v\mapsto L_2v-c_UTv^2$ decreases for $v\ge L_2/(2c_UT)$, and
$\rho_1/T$ lies in this range. So
\begin{equation}\label{eq:GTfar}
G_T(\alpha)\le\frac{L_2\rho_1-c_U\rho_1^2+L_1}T\le\frac{c_F-1}T
\qquad(\alpha\in U,\ |x|\ge\rho_1/T).
\end{equation}
Next let $|x|\le\rho_1/T$, and put $y=x-\Sigma\gamma/T$. Then $\alpha\in U$,
and the bounds on $U$ give
$G_T(\alpha)=\gamma\cdot x-\frac T2x^\top\Sigma^{-1}x+c^*/T+\epsilon$ with
$|\epsilon|\le L_1|x|^2+L_1T|x|^3+L_1|x|/T\le L_1(\rho_1^2+\rho_1^3+\rho_1)T^{-2}$.
Completing the square gives
$\gamma\cdot x-\frac T2x^\top\Sigma^{-1}x
=\frac1{2T}\gamma^\top\Sigma\gamma-\frac T2y^\top\Sigma^{-1}y$. So
\begin{equation}\label{eq:GTnear}
G_T(\alpha)=\frac{c_F}T-\frac T2y^\top\Sigma^{-1}y+O\Bigl(\frac1{T^2}\Bigr)
\qquad(|x|\le\rho_1/T),
\end{equation}
uniformly in $x$. Since $\Sigma^{-1}$ is positive definite,
\eqref{eq:GTfar} and~\eqref{eq:GTnear} give
$G_T\le c_F/T+O(T^{-2})$ on $U$.

\emph{Part \textup{(ii)}, the upper bound.} The part of the proof of
Theorem~\mref{thm:facemax} that treats compositions outside $U$ uses only that
$U$ is a closed ball around $\alpha^*_F$ inside $\Delta_F^\circ$. So there is
$\eta>0$ with $P_N(n)\le X\cdot O(T^{(k-1)/2}e^{-2\eta T})$ for $\supp n=F$ and
$n/N\notin U$. After increasing $T_0$ this is at most
$Xe^{-\eta T}\le Xe^{(c_F-1)/T}$. For $n/N\in U$ we use~\eqref{eq:GT} and the
bound on $G_T$. Hence $\log(M_F/X)\le c_F/T+O(T^{-2}+T^2/N)$.

\emph{Part \textup{(ii)}, the lower bound.} Let $\alpha^0$ be the point with
$x=\Sigma\gamma/T$. Let $n^0_i$ be an integer nearest to $N\alpha^0_i$ for
$i\ge2$, and $n^0_1=N-\sum_{i\ge2}n^0_i$. Then
$|n^0_i-N\alpha^0_i|\le k$ for every $i$. Since $\alpha^0\to\alpha^*_F$ as
$T\to\infty$ and every $\alpha^*_{F,i}$ is positive, $n^0$ has support $F$
after increasing $T_0$ and $N_0$. For $\alpha=n^0/N$ we have $|y|\le k/N$ and
$|x|\le|\Sigma\gamma|/T+k/N\le\rho_1/T$, because $T\le N$. So
\eqref{eq:GTnear} gives $G_T(n^0/N)\ge c_F/T-O(T/N^2+T^{-2})$,
and~\eqref{eq:GT} gives
$\log(M_F/X)\ge\log(P_N(n^0)/X)\ge c_F/T-O(T^{-2}+T^2/N)$. This bound is sufficient for the claimed remainder.

The two bounds give~\meqref{eq:logmaxnext} for $N\ge N_0$ and
$T_0\le T\le N^{1/3}$. For $1\le T\le T_0$ it follows from~\meqref{eq:logmax},
since $c_F/T$ and $1/T$ are $O(T^{-2})$ there, with a constant that depends
on $T_0$. This proves (ii).

\emph{Part \textup{(iii)}.} Now $F$ and $G$ are as in
Theorem~\mref{thm:crossing}, and $k_G=1$ is allowed. Let $T_N\in[1,N^{1/3}]$ be
a zero of $D_N$. The proof of Theorem~\mref{thm:crossing}(ii) shows that
$\frac\Delta{2\delta}L<T_N\le\frac{2\Delta}\delta L$ for large $N$.
Subtracting~\meqref{eq:logmaxnext} for $G$ from~\meqref{eq:logmaxnext} for $F$
at $T=T_N$ gives
\[
0=-\Delta L+\frac\Delta2\log T_N+\delta T_N+\log\frac{K_F}{K_G}
+\frac{c_F-c_G}{T_N}+O\Bigl(\frac1{T_N^2}+\frac{T_N^2}N\Bigr).
\]
We multiply by $T_N$. The error becomes $O(1/T_N+T_N^3/N)=O(1/L)$, which
is~\meqref{eq:crossingnext}.

\emph{Part \textup{(iv)}, a symmetric pair.} Let $F=\{i,j\}$ with
$q_{ij}=q_{ji}=a>0$, $q_i=q_j=q$ and $\mu_i=\mu_j>0$. Then $F$ satisfies (H).
We first prove~\meqref{eq:pairnext}, without the Laplace lemma. Let
$1\le T\le N^{1/3}$ and let $N$ be so large that $qh\le\frac12$ and $I+hQ$ is
stochastic. By Proposition~\mref{prop:pair}, $M_F=\mu_i(1-qh)^{N-1}S_N(u)$ with
$u=ah/(1-qh)$. Put $z=Nu=aT/(1-qh)$. Then $aT\le z\le2aT$ and
$0\le z-aT\le2aqT^2/N$. By~\meqref{eq:planar-bessel},
\[
\log S_N(u)=\log\frac2N+\log\psi(z)+O\Bigl(\frac{T^2}N\Bigr),
\]
since $z(z+2)/(2N)=O(T^2/N)$ and this is at most $\frac12$ for large $N$. By
\cite[Eq.~10.40.1]{dlmf}, as $z\to\infty$,
$I_0(z)=e^z(2\pi z)^{-1/2}(1+\frac1{8z}+O(z^{-2}))$ and
$I_1(z)=e^z(2\pi z)^{-1/2}(1-\frac3{8z}+O(z^{-2}))$. So
$I_0(z)+I_1(z)=2e^z(2\pi z)^{-1/2}(1-\frac1{8z}+O(z^{-2}))$ and
\[
\log\psi(z)=z+\tfrac12\log z+\tfrac12\log\frac2\pi-\frac1{8z}+\rho(z),
\qquad|\rho(z)|\le C_az^{-2}\quad\text{for }z\ge a .
\]
The bound on $\rho$ holds for all $z\ge a$, and not only for large $z$,
because $\rho(z)=O(z^{-2})$ as $z\to\infty$ and $\rho$ is continuous on $[a,\infty)$. The function
$z+\frac12\log z-\frac1{8z}$ has a bounded derivative on $[a,\infty)$. So
replacing $z$ by $aT$ in it changes it by $O(T^2/N)$, and
\[
\log\psi(z)=aT+\tfrac12\log T+\tfrac12\log\frac{2a}\pi-\frac1{8aT}
+O\Bigl(\frac1{T^2}+\frac{T^2}N\Bigr).
\]
Also $(N-1)\log(1-qh)=-qT+O(T^2/N)$. Adding the three parts,
\[
\log M_F=-L+\tfrac12\log T+\log\Bigl(2\mu_i\sqrt{\frac{2a}\pi}\Bigr)-(q-a)T
-\frac1{8aT}+O\Bigl(\frac1{T^2}+\frac{T^2}N\Bigr).
\]
By Corollary~\mref{cor:symface} with $k=2$, $\mu(F)=2\mu_i$ and
$\tree(q_F)=a$, we have $\lambda_F=q-a$ and
$K_F=2\mu_i\cdot2^{3/2}\sqrt a\,(4\pi)^{-1/2}=2\mu_i\sqrt{2a/\pi}$. This
is~\meqref{eq:pairnext}. Now we compare it with~\meqref{eq:logmaxnext} at
$T=L$. Both expand $\log M_F$, so their difference
$(c_F+\frac1{8a})/L$ is $O(L^{-2}+L^2/N)=O(L^{-2})$. Letting
$N\to\infty$ gives $c_F=-1/(8a)$.

As a check we compute the same value from~\meqref{eq:c-alpha}. Number the
states so that $i$ is state $1$. Then $A'=0$, so in the notation of
Appendix~\mref{app:torus} we have $g(z_2)=a^2/z_2$ and
$b(z_2)=\mu_i(1+a/z_2)^2$. At $\alpha=\alpha^*_F=(\frac12,\frac12)$ the radius
is $\theta_2=a$, so $\Psi=a\cos\varphi$ and
$\chi=\mu_ia(1+e^{-\mathrm i\varphi})^2e^{\mathrm i\varphi}
=2\mu_ia(1+\cos\varphi)$. Hence $\mathcal J=a$, $\Psi_{111}=0$,
$\Psi_{1111}=a$, $\chi(\alpha^*_F,0)=4\mu_ia$, $\chi_1=0$ and
$\chi_{11}=-2\mu_ia$, and~\meqref{eq:laplace-c1} gives
$c_1=-\mu_i+\frac12\mu_i=-\frac12\mu_i$ and $c(\alpha^*_F)=-1/(8a)$. Also
$\gamma_F=0$, because $C_F$ is unchanged when the two states are swapped
(Proposition~\mref{prop:constant}(b) with $\nu=\ones$). So the two values
agree.

\emph{Part \textup{(iv)}, the complete graph.} Let $Q=J-dI$, let $F$ have
$k\ge2$ states and let $\mu_i=\mu_0>0$ for $i\in F$. Then $F$ satisfies (H),
and Corollary~\mref{cor:symface} with $\mu(F)=k\mu_0$ and
$\tree(q_F)=k^{k-2}$ gives $\lambda_F=d-k$, $\alpha^*_F$ uniform on $F$ and
$K_F=\mu_0k^{k+1/2}(4\pi)^{-(k-1)/2}$.

We first show that $\gamma_F=0$, so that $c_F=c(\alpha^*_F)$. By
Example~\mref{ex:complete-face},
\[
C_F(\alpha)=\frac{\mu_0\bigl(\sum_{i\in F}\sqrt{\alpha_i}\bigr)^{k/2+1}}
{(4\pi)^{(k-1)/2}\prod_{i\in F}\alpha_i^{3/4}},
\]
a symmetric function of $(\alpha_i)_{i\in F}$. Take $j\ge2$
and swap the states $1$ and $j$. This maps the point
$\alpha^*_F+v(e_j-e_1)$ to $\alpha^*_F-v(e_j-e_1)$, so
$v\mapsto\log C_F(\alpha^*_F+v(e_j-e_1))$ is even and its derivative at $v=0$
vanishes. In the coordinates $\alpha'$ this derivative is the $j$-th entry of
$\gamma_F$.

Now we evaluate~\meqref{eq:c-alpha} at $\alpha^*_F$ by hand. Put $m=k-1$, the dimension of the
torus, and $\varpi=m/k$.

\emph{The integrand.} At $\alpha^*_F$ we have $\theta(\alpha^*_F)=q-\lambda_F\ones$, so every
radius is $\theta_i=(d-1)-(d-k)=m$. In the notation of Lemma~\mref{lem:torus}, $A'=J-I$ of
size $m$, $a_1=b_1=\ones$ and $\mu'=\mu_0\ones$. With $\mathcal D=\operatorname{diag}(z')+I$ the
Sherman--Morrison formula gives $X(z')=(\mathcal D-\ones\ones^\top)^{-1}=\mathcal D^{-1}+\mathcal
D^{-1}\ones\ones^\top\mathcal D^{-1}/(1-s)$, where $s=\sum_{i\ge2}1/(1+z_i)$. So
$g(z')=s/(1-s)$ and $b(z')=\mu_0(1+g)^2=\mu_0(1-s)^{-2}$. Put $z_i=me^{\mathrm i\varphi_i}$ and
$\varepsilon_i=e^{\mathrm i\varphi_i}-1$, so that $1+z_i=k(1+\varpi\varepsilon_i)$, and put
$y_i=\varpi\varepsilon_i/(1+\varpi\varepsilon_i)$ and $Y=\sum_{i\ge2}y_i$. Then
$1/(1+z_i)=(1-y_i)/k$, so $1-s=(1+Y)/k$ and $g=k/(1+Y)-1$. Since $(m^2-1)/k=m-1$,
the functions of~\meqref{eq:A-fT} are
\[
\Psi(\alpha^*_F,\varphi)=\frac1{1+Y}+m-1+\varpi\sum_{i\ge2}\varepsilon_i,\qquad
\chi(\alpha^*_F,\varphi)=\mu_0k^2m^m\,(1+Y)^{-2}e^{\mathrm i\sum_{i\ge2}\varphi_i}.
\]

\emph{Taylor coefficients.} Write $\mathcal S_j=\sum_{i\ge2}\varphi_i^j$. Expanding
$\varepsilon_i=\mathrm i\varphi_i-\frac12\varphi_i^2-\frac{\mathrm i}6\varphi_i^3+\frac1{24}\varphi_i^4+O(\varphi_i^5)$ in
$y_i=\varpi\varepsilon_i-\varpi^2\varepsilon_i^2+\varpi^3\varepsilon_i^3-\varpi^4\varepsilon_i^4+O(\varepsilon_i^5)$
gives
\[
y_i=\mathrm i\varpi\varphi_i+\varpi(\varpi-\tfrac12)\varphi_i^2
+\mathrm i\varpi(\varpi-\varpi^2-\tfrac16)\varphi_i^3
+\varpi(\tfrac1{24}-\tfrac7{12}\varpi+\tfrac32\varpi^2-\varpi^3)\varphi_i^4+O(\varphi_i^5).
\]
We put this into $1/(1+Y)=1-Y+Y^2-Y^3+Y^4-\cdots$ and into
$\log(\chi/\chi(\alpha^*_F,0))=-2\log(1+Y)+\mathrm i\mathcal S_1$. The linear term of $\Psi$ cancels, and
with $\mathcal P_j$ the part of degree $j$ of the Taylor polynomial of $\Psi(\alpha^*_F,\cdot)$ and $\mathcal L_j$
that of $\log(\chi/\chi(\alpha^*_F,0))$,
\begin{align*}
\mathcal P_2&=-\varpi^2(\mathcal S_2+\mathcal S_1^2),\\
\mathcal P_3&=\mathrm i\bigl[\varpi^2(\varpi-1)\mathcal S_3+\varpi^2(2\varpi-1)\mathcal S_1\mathcal S_2+\varpi^3\mathcal S_1^3\bigr],\\
\mathcal P_4&=\varpi^2\bigl[(\tfrac7{12}-\tfrac32\varpi+\varpi^2)\mathcal S_4+(\varpi-\tfrac12)^2\mathcal S_2^2
+(\tfrac13-2\varpi+2\varpi^2)\mathcal S_1\mathcal S_3\\
&\qquad+3\varpi(\varpi-\tfrac12)\mathcal S_1^2\mathcal S_2+\varpi^2\mathcal S_1^4\bigr],\\
\mathcal L_1&=\mathrm i(1-2\varpi)\mathcal S_1,\qquad \mathcal L_2=\varpi(1-2\varpi)\mathcal S_2-\varpi^2\mathcal S_1^2 .
\end{align*}

\emph{The formula as a Gaussian mean.} By the proof of Lemma~\mref{lem:laplace2},
$c_1=\E\mathcal Q_2$ with $\mathcal Q_2=a_0(\mathcal P_4+\frac12\mathcal P_3^2)+A_1\mathcal P_3+A_2$, where $a=\chi$,
$A_j$ is the part of degree $j$ of its Taylor polynomial, and the variable is a centered Gaussian
vector $x$ with covariance $\mathcal J^{-1}$. Since $A_1=a_0\mathcal L_1$ and $A_2=a_0(\mathcal L_2+\frac12\mathcal L_1^2)$,
\[
c(\alpha^*_F)=\frac{c_1}{a_0}=\E\Bigl[\mathcal P_4+\mathcal L_2+\tfrac12(\mathcal P_3+\mathcal L_1)^2\Bigr](x).
\]
Here $\mathcal P_2=-\frac12x^\top\mathcal Jx$ gives $\mathcal J=2\varpi^2(I+\ones\ones^\top)$, of size
$m$, and $\mathcal J^{-1}=\sigma_0(I-\ones\ones^\top/k)$ with $\sigma_0=1/(2\varpi^2)$.

\emph{An exchangeable Gaussian vector.} Let $G_1,\dots,G_k$ be independent standard Gaussian
variables, $\bar G$ their mean and $w_i=G_i-\bar G$. Then $\sum_{i=1}^kw_i=0$ and
$\E[w_iw_j]=\delta_{ij}-1/k$, so $(x_2,\dots,x_k)$ has the law of
$\sqrt{\sigma_0}(w_2,\dots,w_k)$. Put $W_j=\sum_{i=1}^kw_i^j$. Then $\mathcal S_1=-\sqrt{\sigma_0}\,w_1$ and
$\mathcal S_j=\sigma_0^{j/2}(W_j-w_1^j)$ for $j\ge2$. We put these into the Taylor coefficients and use
$\sigma_0\varpi^2=\frac12$, $1-\varpi=1/k$ and $2\varpi-1=(k-2)/k$. The terms in $w_1^3$ cancel in
$\mathcal P_3+\mathcal L_1$ and the terms in $w_1^4$ cancel in $\mathcal P_4$, and
\begin{align*}
\mathcal P_3+\mathcal L_1&=-\frac{\mathrm i\sqrt{\sigma_0}}{2k}\bigl[W_3+(k-2)w_1(W_2-2)\bigr],\\
\mathcal P_4&=\frac1{4\varpi^2}\Bigl[(\tfrac7{12}-\tfrac32\varpi+\varpi^2)W_4+(\varpi-\tfrac12)^2W_2^2\\
&\qquad+(\varpi+1)(\varpi-\tfrac12)w_1^2W_2-(\tfrac13-2\varpi+2\varpi^2)w_1W_3\Bigr],\\
\mathcal L_2&=\frac{1-2\varpi}{2\varpi}(W_2-w_1^2)-\frac12w_1^2 .
\end{align*}

\emph{The moments.} By exchangeability and $\sum_iw_i=0$, every symmetric function $\Upsilon$ of
$w$ has $\E[w_1\Upsilon(w)]=\frac1k\E[\Upsilon(w)\sum_iw_i]=0$. So $\E[w_1W_3]=0$ and
$\E[w_1W_3W_2]=0$. The vector $w$ is a standard Gaussian vector in the hyperplane
$\sum_iw_i=0$, so $W_2$ has the $\chi^2$ law with $m$ degrees of freedom and is independent of
$w/\sqrt{W_2}$. Hence $\E[W_2^j]=m(m+2)\cdots(m+2j-2)$ and
$\E[w_1^2W_2^j]=\E[W_2^{j+1}]/k$, since $\E[w_1^2\mid W_2]=W_2/k$ by symmetry. By Wick's
formula, $\E[w_1^4]=3\varpi^2$, $\E[w_1^6]=15\varpi^3$ and
$\E[w_1^3w_2^3]=9\varpi^2(-\frac1k)+6(-\frac1k)^3$, so $\E[W_4]=3k\varpi^2$ and
$\E[W_3^2]=15k\varpi^3-k(k-1)(9\varpi^2/k+6/k^3)=6m(m-1)/k$. With these values,
\begin{align*}
\E[\mathcal P_4]&=\frac{3k}4\Bigl(\frac7{12}-\frac32\varpi+\varpi^2\Bigr)
+\frac{m(m+2)}{4\varpi^2}\Bigl(\varpi-\frac12\Bigr)\Bigl(\varpi-\frac12+\frac{\varpi+1}k\Bigr)\\
&=\frac{m^4+6m^3-m^2-2}{16mk},\\
\E[\mathcal L_2]&=\frac{1-2\varpi}{2\varpi}\,m\varpi-\frac\varpi2=-\frac{m^2}{2k},\\
\tfrac12\E[(\mathcal P_3+\mathcal L_1)^2]&=-\frac{\sigma_0}{8k^2}\Bigl(\E[W_3^2]+(k-2)^2\E\bigl[w_1^2(W_2-2)^2\bigr]\Bigr)\\
&=-\frac{\sigma_0}{8k^2}\cdot\frac{m(m-1)}k\bigl(m^3+m^2+2m+2\bigr)
=-\frac{m^4+m^2-2}{16mk},
\end{align*}
where we used $\E[w_1^2(W_2-2)^2]=m(m^2+2m+4)/k$ and $\sigma_0=k^2/(2m^2)$. Adding the three,
\[
c_F=c(\alpha^*_F)=\frac{(m^4+6m^3-m^2-2)-8m^3-(m^4+m^2-2)}{16mk}=-\frac{2m^2(m+1)}{16mk}
=-\frac{k-1}8 .
\]
For $k=2$ this is the value $-\frac18$ of the pair. \qed

\begin{thebibliography}{9}

\bibitem{dlmf}
F.~W.~J.~Olver, A.~B.~Olde Daalhuis, D.~W.~Lozier, B.~I.~Schneider,
R.~F.~Boisvert, C.~W.~Clark, B.~R.~Miller, B.~V.~Saunders, H.~S.~Cohl and
M.~A.~McClain (eds.), \emph{NIST Digital Library of Mathematical Functions},
\url{https://dlmf.nist.gov/}.

\end{thebibliography}

\end{document}
